% =====================================================================
% MAIN TEXT -- replaces the current second-domain subsection
% =====================================================================

\subsubsection{CodiEsp Diagnosis Coding}
\label{sec:codiesp}

We instantiate evidence-to-taxonomy retrieval on CodiEsp-D, the diagnosis
portion of the CodiEsp clinical coding benchmark. The official CodiEsp test
cases are the source contexts, and the candidate inventory contains
\num{71344} local diagnosis-code labels after intersecting the CodiEsp valid
diagnosis-code list with the billable FY2018 diagnosis-code inventory. A fact
is $x_i=(X_i,\ell_i,v_i)$: English clinical context $X_i$, a relocated
diagnosis mention $\ell_i$, and target diagnosis code $v_i$. CodiEsp's original
offsets are Spanish offsets, so we first relocate each gold diagnosis mention
into the English context and then study grounding and retrieval given that
located mention; upstream extraction errors therefore do not propagate into any
number we report.

\paragraph{Dimensions and their match types.}
We instantiate the schema of Section~\ref{sec:factorized} with $M{=}6$
dimensions drawn from the diagnosis-code inventory and CodiEsp clinical
evidence (Table~\ref{tab:codiesp_schema}). What matters downstream is not the
dimension names but how each one can be checked against a candidate.
\textsc{Family}, \textsc{Role}, \textsc{Qualifier}, \textsc{Scope}, and
\textsc{Temporal} are matched exactly after normalization against controlled
vocabularies of 21, 9, 21, 5, and 12 categories. \textsc{Event} has no
controlled vocabulary and falls back to token overlap against the candidate's
label and documentation text. These per-dimension match types, rather than the
dimensions themselves, are what a new inventory has to supply.

\paragraph{Category profiles.}
The five vocabulary-matched dimensions give each code a \emph{category
profile}: the tuple of its matched values. Two codes sharing a profile differ
primarily in specific condition wording, whereas codes with different profiles
can differ in chapter, code class, modifier, laterality, or encounter/extension
status. This is what the verifier window of Section~\ref{sec:llmverifier} is
filled by, so that the window spans competing clinical readings instead of only
competing spellings.

Appendix~\ref{app:codiesp} gives the split statistics, the context
serialization, the index and its diagnosis-code pre-filter, the vocabularies,
and the evaluation convention.

\begin{table}[!htbp]
\centering
\scriptsize
\begin{tabular}{@{}llll@{}}
\toprule
\textbf{Dimension} & \textbf{Meaning} & \textbf{Example} & \textbf{Match} \\
\midrule
\textsc{Family} & broad clinical family & Respiratory diseases & vocab (21) \\
\textsc{Role} & diagnosis-code class & Injury/poisoning & vocab (9) \\
\textsc{Event} & specific condition or finding & Costal fracture & overlap \\
\textsc{Qualifier} & modifier or severity & Acute, malignant, open & vocab (21) \\
\textsc{Scope} & laterality & Left, bilateral & vocab (5) \\
\textsc{Temporal} & encounter or extension status & Initial encounter & vocab (12) \\
\bottomrule
\end{tabular}
\caption{The $M{=}6$ dimensions instantiated for CodiEsp diagnosis coding.
\textbf{Match} is how the dimension is compared against a candidate; the five
vocabulary-matched dimensions form the category profile.}
\label{tab:codiesp_schema}
\end{table}


% =====================================================================
% APPENDIX -- replaces/adds the second-domain task details
% =====================================================================

\section{Task Instantiation Details}
\label{app:instantiation}

\subsection{CodiEsp Diagnosis Coding}
\label{app:codiesp}

\paragraph{Splits and statistics.}
The evaluation set is a deterministic exact-relocation slice of the official
CodiEsp test split, not a new random split. We keep all 250 test clinical
cases. From the diagnosis annotations, after inventory filtering,
deduplication, and exact English mention relocation, the split contains
\num{3144} target facts over \num{958} distinct gold codes
(Table~\ref{tab:codiesp_data_stats}). The full prepared diagnosis set contained
\num{3431} facts; \num{287} were dropped because they did not pass the exact
English relocation criterion. No development split or CodiEsp-specific
configuration selection is used.

\begin{table}[!htbp]
\centering
\small
\begin{tabular}{@{}lr@{}}
\toprule
\textbf{Statistic} & \textbf{CodiEsp-D} \\
\midrule
Source documents & 250 \\
Target facts & 3,144 \\
Unique gold codes & 958 \\
Facts per document & 12.58 \\
3-character gold codes & 410 \\
4--5-character gold codes & 2,547 \\
6--7-character gold codes & 187 \\
\bottomrule
\end{tabular}
\caption{CodiEsp diagnosis test statistics after exact English mention
relocation.}
\label{tab:codiesp_data_stats}
\end{table}

\paragraph{Locus and context serialization.}
The locus $\ell_i$ carries the relocated English diagnosis mention and the
local English clinical locus. Because CodiEsp offsets are offsets into Spanish
clinical text, each gold mention is first re-located in the English
machine-translated context with one model call and an exact-substring check.
The retained split has \texttt{relocation\_parse\_ok\_rate}=1.0 and
\texttt{relocation\_exact\_substring\_rate}=1.0; among retained facts,
\num{904} relocations were found from an aligned-sentence candidate scope and
\num{2240} from a document-level candidate scope before exact substring
selection. The source context $X_i$ is serialized as clinical mention, code
class \texttt{diagnosis}, input type, and English source context, with the
runner's default \texttt{12000}-character context budget for query generation
and reranking. Every method receives the same serialization, and direct
retrieval uses it verbatim as its query, so this choice sets the floor as well
as the input to grounding.

\paragraph{Taxonomy filtering and index.}
Of the local diagnosis-code labels available to the CodiEsp setup, we retain
\num{71344} candidates after intersecting the CodiEsp valid diagnosis-code list
with billable FY2018 diagnosis codes. Each code is one document holding its
identifier, canonical English label, documentation text, and structural
metadata. Since the inventory has no single official definition paragraph per
code, the documentation string concatenates, in order, current-code inclusion
terms, inherited include notes, hierarchy path text, and alphabetic-index
lead/sub-term paths resolving to the code. Exclusion notes are kept separate
and are not included in retrieval text. Retrieval is BM25 at depth $K{=}200$,
implemented by the local runner with $k_1=1.5$ and $b=0.75$.

\paragraph{Datatype pre-filter.}
The analogue of the financial datatype pre-filter is the diagnosis-code
inventory restriction. Every fact in this CodiEsp instantiation has code class
\texttt{diagnosis}, and every retained candidate is a billable diagnosis code,
so the shared type-filter hook restricts scoring to that diagnosis inventory.
Where a gold code is not retained by the inventory construction, the fact is
counted as a miss rather than excluded from a method-specific denominator. The
filter is a property of the index and applies identically to every method
compared, direct retrieval included.

\paragraph{Label-coverage term.}
Sparse scoring over compositional diagnosis labels can favor longer labels that
repeat query terms without matching the intended code reading. The retriever
therefore supports adding to the range-normalized BM25 score two token-coverage
terms between the query and the candidate's canonical label, one normalized by
the label's length and one by the query's, under a single weight
$w_{\mathrm{cov}}$. In this CodiEsp matrix we run both
$w_{\mathrm{cov}}=0.0$ and $w_{\mathrm{cov}}=1.0$ cells. Like the diagnosis
pre-filter, the term is a property of the index rather than an FHS component
and is enabled identically for every method within a cell.

\paragraph{Vocabularies and the match operator.}
The five controlled vocabularies were fixed before evaluation: 21
\textsc{Family} categories, 9 \textsc{Role} classes, 21 \textsc{Qualifier}
modifiers, 5 \textsc{Scope} laterality values, and 12 \textsc{Temporal}
encounter/extension values. \textsc{Family} corresponds to broad chapter-level
clinical families; \textsc{Role} distinguishes classes such as
disease/disorder, neoplasm, injury/poisoning, external cause, and
health-status factor; \textsc{Scope} contains right, left, bilateral,
unspecified-side, and not-applicable; and \textsc{Temporal} contains encounter,
sequela, healing, stage, fetus-specific, unspecified, and not-applicable
categories. A normalization map takes the generator's free-form output onto the
controlled values, and values that fail to normalize are logged. For the
vocabulary-free \textsc{Event} dimension, a value matches a candidate when its
tokens overlap the candidate's concatenated label and documentation text.

\paragraph{Evaluation convention.}
A prediction is correct when it matches the gold diagnosis-code identifier
exactly after the runner's tag normalization. Each fact has exactly one gold
code. Recall@10, Recall@50, Recall@200, and MRR are computed over the retrieved
list at depth $K{=}200$ before the shared listwise selector. Accuracy is the
top-1 output of the selector.

\paragraph{Shared backbone and selector.}
All grounding methods use \textsc{Qwen3-32B} as the generator and, where
applicable, the verifier backbone. Downstream of every method with reranking
enabled, a shared listwise selector using the same \textsc{Qwen3-32B} model
receives that method's top-$K$ list and returns at most twenty ranked codes. It
is not part of any method's contribution; the same selector prompt, candidate
format, and evaluation code are held fixed across arms.
